\documentclass[11pt,a4paper]{article}
\usepackage[margin=2.4cm]{geometry}
\usepackage{amsmath,amssymb,booktabs,hyperref}
\newcommand{\dtau}{\Delta\tau}
\newcommand{\code}[1]{\texttt{\small #1}}
\title{Fast integrators with a vertically implicit step\\
on a terrain-following grid:\\ Verlet EI and the WRF forward--backward scheme}
\author{MultiRate.jl}
\date{\today}
\emergencystretch=3em
\begin{document}
\maketitle

\begin{abstract}
We formulate the linearly implicit Verlet step ``Version EI'' of the note \emph{Split-explicit
with implicit treatment of part of the fast waves} (Imex.pdf, 2024) for the fast part of the
Gal-Chen finite-volume discretization, derive the tridiagonal system for the vertical momentum $W$,
discuss when the contravariant vertical mass flux $\Omega$ has to be the unknown instead, and compare
the scheme with the forward--backward acoustic step of WRF with its vertically implicit,
off-centred Crank--Nicolson step. Both are implemented as fast integrators of the frozen MIS
split with the $W$ system and an explicit metric term (Section~\ref{sec:impl}) and compared on the
Agnesi mountain wave (Section~\ref{sec:agnesi}).
\end{abstract}

\tableofcontents

%======================================================================
\section{Introduction}
Split-explicit methods (MIS, RK3) integrate the fast acoustic--gravity part with small steps
$\dtau$. On grids with $\Delta z\ll\Delta x$ the vertical acoustic Courant number
$c_s\dtau/\Delta z$ is far above one, and the vertical part has to be implicit, while the horizontal
part stays explicit: HEVI. Three ingredients decide the properties of such a small step:
\begin{enumerate}
\item the explicit horizontal integrator (Störmer--Verlet, forward--backward);
\item the implicit vertical integrator (Crank--Nicolson, off-centred) and the variable whose
column system is solved ($W$, $\Omega$);
\item the place of the terrain metric terms (explicit or implicit).
\end{enumerate}
The scheme of the Imex note combines the first two in one Verlet step (Section~\ref{sec:ei});
WRF (Klemp, Skamarock and Dudhia 2007; Skamarock et al.\ 2019) uses a forward--backward step with
an off-centred vertically implicit part (Section~\ref{sec:wrf}). The comparison is in
Section~\ref{sec:cmp}.

%======================================================================
\section{The fast system on the terrain-following grid}\label{sec:sys}
Gal-Chen coordinate $z=z_s(x)+\eta J(x)$, $J=1-z_s/H$, C-grid with $\rho,\Theta$ (and $\pi$) at the cell
centres, $U=\rho u$ at the $x$-faces, $W=\rho w$ at the $\eta$-faces; $\Omega=W-z_x\bar U$ is the
contravariant mass flux through the $\eta$-surfaces ($\bar U$ the four-point mean; $\Omega=0$ at the ground
and the top). Thermodynamic variables $t=(\rho,\Theta)$, momenta $(U,W)$. The fast part is
\begin{align*}
\rho_t &= -\tfrac1J\big[\delta_x(J_xU)+\delta_\eta\Omega\big],\\
\Theta_t &= -\tfrac1J\big[\delta_x(\theta_xJ_xU)+\delta_\eta(\theta_\eta\Omega)\big],\\
U_t &= -\tfrac1{J_x}\big[\delta_x(Jp)-\delta_\eta(z_x\,p_\sigma)\big],\\
W_t &= -\frac{\delta_\eta p}{J\Delta\eta}-g\bar\rho ,
\end{align*}
$p_\sigma$ the mean at the $\eta$-faces, $\theta_x,\theta_\eta$ frozen face values, and $p=p(\Theta)$ (full
pressure), $p\approx P\,\Theta$ with frozen $P=(\gamma p/\Theta)_f$ (product form) or the Exner pressure as an
additional variable. We use the blocks
\[
t_t = D_h U + D_v W + f_t,\qquad U_t = G_x\,p + f_U,\qquad W_t = G_z\,p - g\bar\rho + f_W,
\]
$D_h$ the horizontal divergence \emph{including} the part $-z_x\bar U$ of $\Omega$, $D_v$ the part with $W$,
$G_x$ the horizontal pressure gradient \emph{including} the metric term $\delta_\eta(z_xp_\sigma)/J_x$, $G_z$
the vertical one; $f$ the constant forcing of the stage. Two observations shape everything that
follows:
\begin{itemize}
\item a vertical pressure difference acts on $W$ \emph{and}, through the metric term, on $U$, and
$U$ feeds back into the vertical divergence through $\Omega$;
\item for $z_x=0$ the vertical part $(D_v,G_z,g)$ is a closed column system.
\end{itemize}

%======================================================================
\section{Verlet EI}\label{sec:ei}
\subsection{The step}
In the Imex note the momentum operator is split into an explicit part $M^E_{up}$ and an implicit part
$M^I_{up},M^I_{u\rho}$; the divergence $M_{pu}$ is not split. Here $M_{pu}=D_h+D_v$, $M^E_{up}=G_x$ ($U$ rows),
$M^I_{up}=G_z$, $M^I_{u\rho}=-g\,\overline{(\cdot)}$ ($W$ rows). With $h=\dtau$ and the off-centring $\gamma$
(Version III; $\gamma=\tfrac12$ is Version EI):
\begin{align}
t^{n+1/2} &= t^n+\tfrac h2\big[D_hU^n+D_vW^n+f_t\big],\label{eq:ei1}\\
U^{n+1} &= U^n+h\big[G_x\,p(t^{n+1/2})+f_U\big],\label{eq:ei2}\\
W^{n+1} &= W^n+h\big[(1-\gamma)(G_zp^n-g\bar\rho^n)+\gamma(G_zp^{n+1}-g\bar\rho^{n+1})\big]+hf_W,\label{eq:ei3}\\
t^{n+1} &= t^{n+1/2}+\tfrac h2\big[D_hU^{n+1}+D_vW^{n+1}+f_t\big].\label{eq:ei4}
\end{align}
The thermodynamic half steps are explicit in the full divergence; only $W$ is implicit.

\subsection{The system for $W$}
$U^{n+1}$ is known after \eqref{eq:ei2}. With $t^*=t^{n+1/2}+\tfrac h2[D_hU^{n+1}+f_t]$, \eqref{eq:ei4} reads
$t^{n+1}=t^*+\tfrac h2D_vW^{n+1}$. Linearize $p^{n+1}=p(t^*)+P\,(\Theta^{n+1}-\Theta^*)$ (exact for the product
form and the Exner form; for the full pressure linearly implicit with $P=\gamma p/\Theta$ at the start of the
stage) and insert into \eqref{eq:ei3}:
\[
\Big[I-\tfrac{\gamma h^2}{2}\big(G_zP\,D_v^\Theta-g\,\overline{D_v^\rho}\big)\Big]W^{n+1}
=W^n+h(1-\gamma)\big(G_zp^n-g\bar\rho^n\big)+h\gamma\big(G_zp(t^*)-g\bar\rho^*\big)+hf_W .
\]
With $D_v^\rho W=-\delta_\eta W/(J\Delta\eta)$, $D_v^\Theta W=-\delta_\eta(\theta_\eta W)/(J\Delta\eta)$ and the face pressure
difference $\alpha_k\Theta_k-\beta_k\Theta_{k-1}$ ($\alpha=\beta=P_f$ for the product form) this is, at the face $k$
of a column with $\Delta z=J\Delta\eta$,
\[
W_{k-1}\big[-r\beta_k\theta_{k-1}+g_g\big]+W_k\big[1+r\theta_k(\alpha_k+\beta_k)\big]+W_{k+1}\big[-r\alpha_k\theta_{k+1}-g_g\big]=b_k,
\]
\[
r=\frac{\gamma h^2}{2\Delta z^2},\qquad g_g=\frac{\gamma h^2g}{4\Delta z},
\]
a tridiagonal system per column; $\rho^{n+1},\Theta^{n+1}$ follow from $t^{n+1}=t^*+\tfrac h2D_vW^{n+1}$.
Equivalently: with the row scaling $D=\mathrm{diag}(\tfrac h2\ \text{on }t,\ \gamma h\ \text{on }W)$ the
steps \eqref{eq:ei3}--\eqref{eq:ei4} are one solve $(I-DJ_v)Y=DG_v$ of the vertical part, preceded by the
explicit vertical increments $\tfrac h2D_vW^n$, $(1-\gamma)hG_v$ in \eqref{eq:ei1}; the $W$ system depends on
the product $\tfrac h2\cdot\gamma h$ only. The elimination in detail, for all forms, is in Section~\ref{sec:col}.

\subsection{Properties}
\begin{itemize}
\item For a purely vertical problem ($D_h=G_x=0$): $t^{n+1}=t^n+\tfrac h2D_v(W^n+W^{n+1})$ and
\eqref{eq:ei3}: with $\gamma=\tfrac12$ this is the trapezoidal rule (Crank--Nicolson) of the column system,
unconditionally stable and non-dissipative.
\item For a purely horizontal problem it is Störmer--Verlet (position Verlet), symmetric, second
order, stable for $|\omega h|\le2$.
\item The step is symmetric for $\gamma=\tfrac12$ and of second order; for $\gamma\ne\tfrac12$ of first order.
\item One column solve and two evaluations of the horizontal thermodynamic part, one of the
horizontal momentum part, two of the vertical part per small step. The present Strang
composition (\code{:hevi}: horizontal Verlet half step, vertical Crank--Nicolson, horizontal half
step) needs one solve, four horizontal thermodynamic and two horizontal momentum evaluations.
\end{itemize}

%======================================================================
\section{$W$ or the contravariant $\Omega$?}\label{sec:omega}
\subsection{Metric term explicit: the two are equivalent}
In \eqref{eq:ei2} the metric term $\delta_\eta(z_xp_\sigma)/J_x$ is part of $G_x$ and explicit. Then
$\bar U^{n+1}$ is known before the column solve, and $\Omega^{n+1}=W^{n+1}-z_x\bar U^{n+1}$ is an affine shift:
the system for $\Omega$ has the same matrix and the right-hand side shifted by the known
$z_x\bar U^{n+1}$. Nothing is gained by solving for $\Omega$.

\subsection{The price of the explicit metric term}
Constant slope $s=z_x$, no gravity, small amplitude, vertical structure only. The fast system
contains
\[
p_t=-c_s^2\rho_0\,\partial_z(W-sU),\qquad W_t=-\frac{\partial_zp}{\rho_0},\qquad U_t=s\,\frac{\partial_zp}{\rho_0}
\quad\Longrightarrow\quad \Omega_t=-(1+s^2)\frac{\partial_zp}{\rho_0},
\]
the vertical acoustic mode of the terrain-following coordinate has the frequency $c_sk_z\sqrt{1+s^2}$,
and the part $s^2$ of the restoring force passes through the explicit $U$. A heuristic bound for the
explicit coupling, from the largest vertical wavenumber $k_z\approx2/(J\Delta\eta)$, is
\[
\frac{c_s\,\dtau\,|z_x|}{J\Delta\eta}\lesssim1 .
\]
(This is an estimate, not yet the stability analysis of the coupled step.) Agnesi case of the
MIS study: $350\cdot4.3\cdot0.026/200\approx0.2$; steep terrain, $z_x=0.5$, $\Delta z=50$\,m:
$\approx15$. For steep terrain the metric term has to be implicit.

\subsection{Metric term implicit: the system for $\Omega$}
Move the vertical pressure difference of the $U$ equation into the implicit part,
\[
M^I_{up}=\big(\delta_\eta(z_xp_\sigma)/J_x\ \text{on }U,\ G_z\ \text{on }W\big),\qquad M^E_{up}=-\delta_x(Jp)/J_x .
\]
The
vertical divergence depends on the momenta only through $\Omega$ ($\delta_\eta\Omega$, $\delta_\eta(\theta_\eta\Omega)$), and
\[
\Omega^{n+1}=\Omega'+h\big[(1-\gamma)G_\Omega p^n+\gamma G_\Omega p^{n+1}\big],\qquad
G_\Omega p=G_zp-z_x\,\overline{\tfrac1{J_x}\delta_\eta(z_xp_\sigma)}\ \approx\ -(1+z_x^2)\frac{\delta_\eta p}{J\Delta\eta},
\]
$\Omega'$ containing all explicit contributions. This is a tridiagonal system for $\Omega$ with the metric factor
$1+z_x^2$ on the vertical pressure gradient; $U^{n+1}$ is completed afterwards with the known
$p^{n+1}$. Exactly column-local it is only with two approximations:
\begin{enumerate}
\item[(a)] the four-point means ($\bar U$ in $\Omega$, $p_\sigma$ and the average to the $W$ face in the metric
term) are taken within the column in the implicit part; exactly they couple neighbouring columns;
\item[(b)] in \eqref{eq:ei4} the horizontal divergence $D_hU^{n+1}$ uses $U$ before its implicit metric
increment. With the full $M_{pu}$, as in the Imex note
($\mathcal M=I-\tfrac14\dtau^2M^I_{up}M_{pu}$), the system for $(U,W)$ is two-dimensional.
\end{enumerate}
With (b) the step is no longer exactly symmetric, and the energy argument of the C-grid has to be
checked again.

%======================================================================
\section{The WRF forward--backward scheme}\label{sec:wrf}
\subsection{WRF}
WRF integrates the acoustic modes in the small steps of its RK3 (Wicker and Skamarock 2002), in a
terrain-following \emph{hydrostatic-pressure (mass)} coordinate with column mass $\mu$, geopotential
$\phi$ and perturbation variables about the state of the large step (Klemp, Skamarock and Dudhia
2007; Skamarock et al.\ 2019). Schematically, one small step $\tau\to\tau+\dtau$:
\begin{enumerate}
\item \emph{horizontal momentum forward}: $U^{\tau+\dtau}=U^\tau+\dtau\,[\text{horizontal pressure gradient
of }p^{*}+\text{metric term }\partial_\eta p\,\partial_x\phi+R_U]$, with the divergence damping
$p^*=p^\tau+\gamma_d(p^\tau-p^{\tau-\dtau})$, $\gamma_d=0.1$; the metric term is explicit;
\item \emph{column mass and $\Omega$ backward}: $\mu^{\tau+\dtau}$ from the vertically integrated horizontal
divergence of $U^{\tau+\dtau}$, then $\Omega^{\tau+\dtau}$ diagnostically from the continuity equation of the
column --- in the mass coordinate $\Omega$ is \emph{not} an unknown of the implicit system;
\item $\Theta^{\tau+\dtau}$ with the horizontal fluxes of $U^{\tau+\dtau}$ and the vertical ones of $\Omega^{\tau+\dtau}$;
\item \emph{vertically implicit}: $W$ and $\phi$ with the off-centred average
$\overline a^\tau=\tfrac{1+\beta_d}2a^{\tau+\dtau}+\tfrac{1-\beta_d}2a^\tau$, $\beta_d=0.1$, in the vertical pressure
gradient (pressure from the linearized equation of state in $\phi_\eta$, $\Theta$, $\mu$), the buoyancy and the
geopotential equation $\phi_t\propto gW-\Omega\phi_\eta$; eliminating $\phi$ gives a tridiagonal system for $W$.
\end{enumerate}
In addition: an external-mode filter and, on the large step, the RK3 with $n_s/3$, $n_s/2$, $n_s$ small
steps. The vertical unknown is the Cartesian $W$; the terrain enters the implicit system only
through the explicit metric term and the diagnostic $\Omega$.

\subsection{The analogue on the Gal-Chen grid}
In height coordinates $\Omega=W-z_x\bar U$ is not diagnostic; it couples $W$ to the thermodynamic variables.
The analogue (\code{:wrf} in \code{TimeStepping/fast.jl}, linearized fast operator), with $\beta=(1+\beta_d)/2=0.55$:
\begin{align*}
U^{n+1}&=U^n+h\big[G_x\,p(t^{*})+f_U\big],\qquad t^*=t^n+\gamma_d(t^n-t^{n-1}),\\
\begin{pmatrix}t\\W\end{pmatrix}^{n+1}&=\begin{pmatrix}t\\W\end{pmatrix}^n
+h\begin{pmatrix}D_hU^{n+1}+f_t\\ f_W\end{pmatrix}
+h\,A_v\Big[(1-\beta)\begin{pmatrix}t\\W\end{pmatrix}^n+\beta\begin{pmatrix}t\\W\end{pmatrix}^{n+1}\Big],
\end{align*}
$A_v$ the vertical part $(D_v;\,G_z,-g)$. Elimination gives the $W$ system of Section~\ref{sec:ei} with
$r=\beta^2h^2/\Delta z^2$, $g_g=\beta^2h^2g/(2\Delta z)$; the metric term is explicit (as in WRF).

%======================================================================
\section{Linear algebra}\label{sec:la}
\subsection{The column system}\label{sec:col}
In a column $i$ with $\Delta z=J_i\Delta\eta$, cells $k=1,\dots,n_z$ and faces $k=1,\dots,n_z+1$ ($W_1=W_{n_z+1}=0$),
the vertical part $G_v$ acts only on the thermodynamic variables $t$ and on $W$; with frozen coefficients
it is linear,
\[
J_v=\begin{pmatrix}0&D_v\\ B_v&0\end{pmatrix}\quad\text{on }\begin{pmatrix}t\\W\end{pmatrix},
\]
\[
(D_vW)_{\phi,k}=-\frac{\phi_{k+1}W_{k+1}-\phi_kW_k}{\Delta z},\qquad
(B_vt)_k=-\frac{\delta p_k}{\Delta z}-\frac g2(\rho_{k-1}+\rho_k),
\]
$\phi$ the frozen face value of the variable ($\phi=1$ for $\rho$, $\theta_\eta$ for $\Theta$, $q_{v,\eta}$, $q_{c,\eta}$ for the moist
variables), and the face pressure difference
\begin{center}\small
\begin{tabular}{lll}\toprule
form & $\delta p_k$ & coefficients\\\midrule
product & $\alpha_k\Theta_k-\beta_k\Theta_{k-1}$ & $\alpha_k=\beta_k=(\partial p/\partial\Theta)_{f,k}$, frozen\\
\code{:gradp} & $p(\Theta_k)-p(\Theta_{k-1})$ & linearized: $\alpha_k\Theta_k-\beta_k\Theta_{k-1}$,\\
 & & $\alpha_k=(\gamma p/\Theta)_k$, $\beta_k=(\gamma p/\Theta)_{k-1}$ at the stage start\\
Exner, moist & $P_k(\pi_k-\pi_{k-1})$ & $P_k=c_p\Theta_{f,k}$, $(D_vW)_{\pi,k}=-(c^B_{k+1}W_{k+1}-c^T_kW_k)/\Delta z$\\\bottomrule
\end{tabular}
\end{center}
For the full pressure the matrix is the Jacobian at the stage start (linearly implicit); for the other forms it is exact.

\paragraph{Two-factor solve.} All schemes need $Y=(I-DJ_v)^{-1}R$ with the row scaling $D=\mathrm{diag}(c_tI_t,\,c_wI_W)$:
\[
\begin{pmatrix}I&-c_tD_v\\-c_wB_v&I\end{pmatrix}\begin{pmatrix}Y_t\\Y_W\end{pmatrix}=\begin{pmatrix}R_t\\R_W\end{pmatrix}.
\]
The first block row gives $Y_t=R_t+c_tD_vY_W$; inserted into the second (Schur complement on the $W$ block):
\[
\big(I-c_tc_w\,B_vD_v\big)\,Y_W=R_W+c_w\,B_vR_t .
\]
With $r=c_tc_w/\Delta z^2$ and $g_g=c_tc_wg/(2\Delta z)$, row $k=2,\dots,n_z$ of the product and \code{:gradp} forms is
\begin{multline*}
\big[-r\beta_k\theta_{k-1}+g_g\big]Y_{W,k-1}+\big[1+r\theta_k(\alpha_k+\beta_k)\big]Y_{W,k}+\big[-r\alpha_k\theta_{k+1}-g_g\big]Y_{W,k+1}\\
=R_{W,k}-\frac{c_w}{\Delta z}\big(\alpha_kR_{\Theta,k}-\beta_kR_{\Theta,k-1}\big)-\frac{c_wg}2\big(R_{\rho,k-1}+R_{\rho,k}\big),
\end{multline*}
of the Exner and moist forms
\begin{multline*}
\big[-rP_kc^T_{k-1}+g_g\big]Y_{W,k-1}+\big[1+rP_k(c^T_k+c^B_k)\big]Y_{W,k}+\big[-rP_kc^B_{k+1}-g_g\big]Y_{W,k+1}\\
=R_{W,k}-\frac{c_w}{\Delta z}P_k\big(R_{\pi,k}-R_{\pi,k-1}\big)-\frac{c_wg}2\big(R_{\rho,k-1}+R_{\rho,k}\big)
\end{multline*}
(the coefficient of $Y_{W,k-1}$ for $k=2$ and of $Y_{W,k+1}$ for $k=n_z$ vanish). The pressure part is a discrete
$-c_s^2\partial_{zz}$ with positive diagonal, the gravity part $\pm g_g$ is antisymmetric and small
($g_g/r=g\Delta z/2\ll c_s^2$); the system is solved by the Thomas algorithm without pivoting, factored once per stage and
value of $c_tc_w$ (\code{\_frozen\_factor}, \code{\_exner\_factor}, cached under the key $c_tc_w$). Back substitution:
\[
Y_{\phi,k}=R_{\phi,k}-\frac{c_t}{\Delta z}\big(\phi_{k+1}Y_{W,k+1}-\phi_kY_{W,k}\big)\quad\text{for every thermodynamic variable }\phi
\]
($\rho$, $\Theta$, $\pi$ with $c^B,c^T$, $q_v$, $q_c$); $U$ rows: $Y_U=R_U$. Moisture enters $B_v$ only through $\rho$ and
$\pi$, so the moist variables are back substitution only. The matrix depends on $c_t$ and $c_w$ only
through the product; the schemes use
\begin{center}\small
\begin{tabular}{lccc}\toprule
& $c_t$ & $c_w$ & $c_tc_w$\\\midrule
HEVI (Crank--Nicolson step of length $h$) & $h/2$ & $h/2$ & $h^2/4$\\
EI & $h/2$ & $\gamma h$ & $\gamma h^2/2$\\
WRF & $\beta h$ & $\beta h$ & $\beta^2h^2$\\\bottomrule
\end{tabular}
\end{center}
For $\gamma=\beta=\tfrac12$ the three matrices are identical.

\subsection{The steps as compositions}
Let $H$ be the linear (affine with the forcing) part ``thermodynamics and vertical'', acting on
$(t,W)$ with $U$ as a parameter,
\[
H(t,U,W)=\big(D_hU+D_vW+f_t,\ 0,\ B_vt+f_W\big),
\]
and define for a row scaling $D=(c_t,c_w)$ the explicit and the implicit step of $H$ and the momentum kick,
\[
E_D\,z=z+D\,H(z),\qquad I_D\,z=z^+\ \text{with}\ z^+=z+D\,H(z^+),\qquad K_h\,z=z+h\,(0,\,G_xp(t)+f_U,\,0).
\]
$I_D$ is one two-factor solve: $z^+=z+(I-DJ)^{-1}DH(z)$; $U$ is constant in $E_D$ and $I_D$, so the horizontal
thermodynamic term $D_hU$ is the same explicit or implicit. Then (Sections~\ref{sec:ei}, \ref{sec:wrf})
\[
\boxed{\ \text{EI}(\gamma)=I_{(h/2,\,\gamma h)}\circ K_h\circ E_{(h/2,\,(1-\gamma)h)},\qquad
\text{WRF}(\beta)=I_{\beta h}\circ E_{(1-\beta)h}\circ K_h\ }
\]
($I_{\beta h}E_{(1-\beta)h}$ is the $\beta$-weighted Crank--Nicolson step of $H$ with the new $U$, the second line of
the WRF step). Without $H$ ($D_v=B_v=0$) EI is position Verlet and WRF forward--backward; without $K$ both
are the (off-centred) Crank--Nicolson step of the column.

\paragraph{Equivalence for $\gamma=\beta=\tfrac12$.} With $c=h/2$ and $E_c=I+cH$, $I_c=(I-cH)^{-1}$ (affine: the
constant parts commute as well, $(I+cH)(I-cH)^{-1}cf+cf=(I-cH)^{-1}2cf$), $E_c$ and $I_c$ commute, and
\[
\text{WRF}(\tfrac12)=I_cE_cK_h=E_cI_cK_h=E_c\,\big(I_cK_hE_c\big)\,E_c^{-1}=E_c\,\text{EI}(\tfrac12)\,E_c^{-1}.
\]
This holds for arbitrary linear operators $D_h,D_v,G_x,B_v$ --- terrain (the $-z_x\bar U$ part of $\Omega$ and the
metric term), gravity, sponge-free forcing --- not only for Fourier modes. Consequences:
\begin{itemize}
\item the two small steps have the same eigenvalues: the same stability limit, the same dispersion
and (none) dissipation of every mode;
\item they differ by the change of variables $E_c$ (an explicit half step of $H$): WRF's thermodynamic
variables are EI's $t^{n+1/2}$, WRF's $W$ is EI's $W+\tfrac h2B_vt$. Over a stage of $n$ small steps
\[
\text{EI}^n=E_c^{-1}\,\text{WRF}^n\,E_c=I_cK_h\,\text{WRF}(\tfrac12)^{n-1}\,E_c ;
\]
\item the combined half steps of Section~\ref{sec:impl} are exactly this formula: $E_c$ at the stage start,
then $n-1$ WRF steps ($E_cI_cK_h$; the update $Z\leftarrow2Z-Z_s$ is $E_c$ applied after $I_c$, since $I_c$ gives
$cH(z^+)=z^+-z_s$), and the last step without $E_c$. EI with combined half steps \emph{is} WRF($\tfrac12$),
started and finished correctly;
\item in a MIS or RK3 stage $t$, $U$, $W$ start at the same time level. EI integrates exactly the fast
flow of these variables to second order; WRF integrates $E_c$-transformed variables but starts from
untransformed ones and returns untransformed ones: the stage map differs from EI's by
$\text{WRF}^n-E_c^{-1}\text{WRF}^nE_c=O(\dtau)$ --- the first-order ``conjugation defect'' of Section~\ref{sec:cmp},
and the only difference between the two schemes.
\end{itemize}
Verified in the code (frozen loop, hill grid, 10 steps of 2\,s with constant forcing):
$\|\text{WRF}^n(E_cx)-E_c(\text{EI}^nx)\|/\|x\|=9\cdot10^{-16}$ (product), $3\cdot10^{-15}$ (Exner, moist),
$2\cdot10^{-6}$ (\code{:gradp}, linearly implicit), while $\|\text{WRF}^nx-\text{EI}^nx\|/\|x\|=4$--$8\cdot10^{-4}$.

\subsection{The options in both variants: off-centring and extrapolation}\label{sec:opt}
Both variants have the same two optional parameters, default zero:
\begin{itemize}
\item the \emph{off-centring} $\varepsilon\ge0$ (WRF's \code{epssm}), $\beta=(1+\varepsilon)/2$;
\item the \emph{explicit extrapolation} $\gamma_d\ge0$ (divergence damping, WRF's \code{smdiv}): the momentum kick
uses the extrapolated thermodynamic state $t+\gamma_d(t-t_{\rm prev})$, $t_{\rm prev}$ the state of the previous kick,
\[
K_h^{\gamma_d}z=z+h\,\big(0,\ G_xp(t+\gamma_d(t-t_{\rm prev}))+f_U,\ 0\big).
\]
\end{itemize}
They enter the two schemes in the same places:
\[
\boxed{\ \text{WRF}(\varepsilon,\gamma_d)=I_{\beta h}\circ E_{(1-\beta)h}\circ K_h^{\gamma_d},\qquad
\text{EI}(\varepsilon,\gamma_d)=I_{\beta h}\circ K_h^{\gamma_d}\circ E_{(1-\beta)h}\ }
\]
with scalar weights on all rows of $H$ (thermodynamic rows including $D_hU$ and $f_t$, $W$ rows including $f_W$).
\begin{itemize}
\item \emph{Off-centring.} In WRF, $I_{\beta h}E_{(1-\beta)h}$ is the $\beta$-weighted Crank--Nicolson step of the column
with the new $U$: $(1-\beta)$ old and $\beta$ new values of $D_vW$ in $t$ and of $B_vt$ in $W$. In EI the first
(explicit) part gets the weight $(1-\beta)h$ and the last (implicit) part $\beta h$, on all rows; $D_hU$ is
constant within $E$ and $I$, so only the split of the horizontal divergence between the old and the new
$U$ changes ($\tfrac h2,\tfrac h2\to(1-\beta)h,\beta h$). For $\varepsilon=0$ EI is Verlet EI (Imex version 1).
\item \emph{Extrapolation.} In WRF the kick uses $\tau^n+\gamma_d(\tau^n-\tau^{n-1})$ with $\tau^{-1}=\tau^0$ (no
extrapolation in the first step of a stage). In EI the kick acts on $t^{n+1/2}$, so the extrapolation is
$t^{n+1/2}+\gamma_d(t^{n+1/2}-t^{n-1/2})$, started at the first $t^{1/2}$ (no extrapolation in the first step).
\end{itemize}
\paragraph{Conjugacy for all $(\varepsilon,\gamma_d)$.} $E_{(1-\beta)h}=I+(1-\beta)hH$ and $I_{\beta h}=(I-\beta hH)^{-1}$ are
functions of the same $H$ and commute; for the affine parts, with $H z=Jz+a$,
\[
E_{(1-\beta)h}I_{\beta h}z-I_{\beta h}E_{(1-\beta)h}z
=\big[(I+(1-\beta)hJ)\beta h+(1-\beta)h(I-\beta hJ)-h\,I\big](I-\beta hJ)^{-1}a=0 .
\]
Hence, with $E=E_{(1-\beta)h}$,
\[
\text{WRF}(\varepsilon,\gamma_d)=E\,\text{EI}(\varepsilon,\gamma_d)\,E^{-1},\qquad \text{EI}^n=E^{-1}\text{WRF}^nE=I_{\beta h}K_h\,\text{WRF}^{n-1}E .
\]
For the extrapolation the history has to correspond as well: if the WRF state is $w^n=Ez^n$ ($z^n$ the EI
state), WRF's $\tau^n$ is EI's $t^{n+1/2}$ for all $n$, and both extrapolations start without extrapolation;
the kicks are identical. So for every choice of the two options EI and WRF have the same small
step up to the change of variables $E$, the same eigenvalues, and the same stability; they differ only
at the stage boundaries, where WRF starts and ends without $E$ ($O(\dtau)$ for $\varepsilon=0$, $O((1-\beta)\dtau)$
in general). For $\varepsilon>0$ both are first order in $\dtau$; for $\varepsilon=0$ EI is second order and WRF first order
(only through the stage boundaries).

Verified in the code (frozen loop, hill grid, 10 steps of 2\,s, constant forcing, combined and separate
half steps of EI): $\|\text{WRF}^n(Ex)-E(\text{EI}^nx)\|/\|x\|\le3\cdot10^{-15}$ for $(\varepsilon,\gamma_d)=(0,0),(0.1,0),(0,0.1),(0.1,0.1),
(0.2,0.3)$ in the product, Exner, pressure and moist forms; $2$--$3\cdot10^{-6}$ for \code{:gradp} (linearly implicit).

\paragraph{The Imex version III ($\gamma$, $W$ rows only).} EI also keeps the off-centring of the Imex note,
\code{EI\_GAMMA[]} $=\gamma$, used only with $\varepsilon=0$: $E_{(h/2,(1-\gamma)h)}$ and $I_{(h/2,\gamma h)}$ with different weights on the
thermodynamic and the $W$ rows. These do not commute for $\gamma\ne\tfrac12$: version III is not conjugate to WRF
for any $\beta$, and damps the vertical modes less (Fourier analysis below).

\paragraph{Combined half steps with the options.} With the weights $(e_t,e_w)$ of $E$ and $(i_t,i_w)$ of $I$ (EI($\varepsilon$):
$e_t=e_w=(1-\beta)h$, $i_t=i_w=\beta h$; version III: $(\tfrac h2,(1-\gamma)h)$, $(\tfrac h2,\gamma h)$), $I$ adds $z^+-z_s=i_tH_t+D_{(i_t,i_w)}G_v(z^+)$
and the next $E$ needs $e_tH_t+D_{(e_t,e_w)}G_v(z^+)$, so between two small steps
\[
z\leftarrow z^++\frac{e_t}{i_t}\,(z^+-z_s)+D_{(0,\ e_w-e_ti_w/i_t)}G_v(z^+),
\]
the last term vanishing for EI($\varepsilon$) and being $(1-2\gamma)hB_v(t^+)$ for version III; $\frac{e_t}{i_t}=\frac{1-\beta}\beta$ ($=1$ for $\varepsilon=0$:
$2z^+-z_s$). The extrapolation needs nothing extra.

\subsection{Fourier analysis}
For the acoustic model $t'=-\mathrm i(k_xu+k_zw)c_s$, $u'=-\mathrm ik_xc_st$, $w'=-\mathrm ik_zc_st$ with $a_x=k_xc_sh$,
$a_z=k_zc_sh$, $\nu=(1-\gamma)+\gamma\lambda$, $\mu=(1-\beta)+\beta\lambda$, the ansatz $z^n=\lambda^nz$ gives, besides the
neutral mode $\lambda=1$,
\[
\text{EI}(\gamma):\ (\lambda-1)^2+a_x^2\lambda+\tfrac{a_z^2}2\,\nu\,(1+\lambda)=0,\qquad
\text{WRF}(\beta):\ (\lambda-1)^2+a_x^2\lambda+a_z^2\,\mu^2=0 .
\]
(EI: from $t^{n+1/2}(\lambda-1+\tfrac{a_x^2}2)=t^n(\lambda-1-\tfrac{a_z^2}2\nu)$ and
$t^n\lambda\big(\lambda-1+\tfrac{a_z^2}2\nu\big)=t^{n+1/2}\big(\lambda-1-\tfrac{a_x^2}2\lambda\big)$; the resultant factors as
$(\lambda-1)[\dots]$.) For $\gamma=\beta=\tfrac12$, $\nu(1+\lambda)/2=\mu^2=(1+\lambda)^2/4$: identical polynomials, as proven above.
Horizontally ($a_z=0$) both are stable for $a_x\le2$; vertically unconditionally. The modulus of the
damped pair is the root of the ratio of constant and leading coefficient, $|\lambda|^2=(1+(1-\gamma)a_z^2/2)/(1+\gamma a_z^2/2)$
(EI) and $(1+(1-\beta)^2a_z^2)/(1+\beta^2a_z^2)$ (WRF) for $a_x=0$ --- the formulas of Section~\ref{sec:cmp}. Spectral radius of
the non-neutral modes:
\begin{center}\small
\begin{tabular}{cccccccc}\toprule
$a_x$&$a_z$& EI $\gamma=0.5$ & EI $0.55$ & EI $0.6$ & WRF $\beta=0.55$ & WRF $0.5$\\\midrule
0.0 & 1.0 & 1.0000 & 0.9802 & 0.9608 & 0.9608 & 1.0000\\
0.0 & 3.0 & 1.0000 & 0.9330 & 0.8699 & 0.8708 & 1.0000\\
0.0 & 20.0 & 1.0000 & 0.9054 & 0.8182 & 0.8198 & 1.0000\\
1.0 & 0.3 & 1.0000 & 0.9978 & 0.9956 & 0.9956 & 1.0000\\
1.5 & 3.0 & 1.0000 & 0.9330 & 0.8699 & 0.8708 & 1.0000\\
1.5 & 20.0 & 1.0000 & 0.9054 & 0.9512 & 0.8198 & 1.0000\\\bottomrule
\end{tabular}
\end{center}
EI with $\gamma\approx0.6$ damps like WRF with $\beta=0.55$, except for strongly coupled modes (large $a_x$ and $a_z$),
where EI's damping is weaker. On a grid (max over $a_z\in[0,200]$) the horizontal limit is $a_x=2$ for all
variants.

%======================================================================
\section{Comparison}\label{sec:cmp}
\subsection{Horizontal: the same spectrum}
For the oscillator $p'=-\omega u$, $u'=\omega p$, $a=\omega h$, forward--backward (momentum first) and Verlet
(thermodynamic half step, momentum, half step) have the matrices
\[
M_{FB}=\begin{pmatrix}1-a^2&-a\\a&1\end{pmatrix},\qquad
M_V=\begin{pmatrix}1-\tfrac{a^2}2&-a\big(1-\tfrac{a^2}4\big)\\a&1-\tfrac{a^2}2\end{pmatrix},
\]
both with trace $2-a^2$ and determinant 1: the same eigenvalues, stable for $|a|\le2$, no damping. They
are conjugate, $M_{FB}=T_{1/2}M_VT_{1/2}^{-1}$ with the thermodynamic half step $T_{1/2}$: forward--backward
is Verlet with the thermodynamic variables shifted by half a step. With the vertical part the same holds
exactly with the transformation $E_c$, an explicit half step of the thermodynamic and vertical part
(Section~\ref{sec:la}): EI($\tfrac12$) and WRF($\tfrac12$) are conjugate for the full linear system. Within a fast stage the
conjugation acts only at the beginning and the end; at the stage boundaries of a MIS or RK3 step
(where $t$ and $U$ start at the same time level) it leaves an error of first order in $\dtau$ in the fast
modes, which does not accumulate over the small steps. Verlet EI has no such defect.

\subsection{Vertical: damping}
For the vertical oscillator ($p'=-\omega W$, $W'=\omega p$, $a=\omega h$), both schemes are unconditionally stable;
the squared modulus of the eigenvalues is the determinant of the step matrix:
\begin{align*}
\text{WRF (}t\text{ and }W\text{ off-centred, }\beta):&\quad |\lambda|^2=\frac{1+(1-\beta)^2a^2}{1+\beta^2a^2},\\
\text{EI (}W\text{ off-centred, }\gamma):&\quad |\lambda|^2=\frac{1+(1-\gamma)\,a^2/2}{1+\gamma\,a^2/2}.
\end{align*}
Both are 1 for $\beta=\gamma=\tfrac12$. For small $a$: $1-(2\beta-1)a^2$ against $1-(2\gamma-1)a^2/2$; for
$a\to\infty$: $(1-\beta)^2/\beta^2$ against $(1-\gamma)/\gamma$. WRF's $\beta=0.55$ gives $|\lambda|\to0.82$, EI with
$\gamma=0.55$ only $0.90$; the same limit needs $\gamma=1/(1+(1-\beta)^2/\beta^2)\approx0.60$. Off-centring makes both
schemes first order in $\dtau$.

\subsection{Summary}
\begin{center}\small
\begin{tabular}{p{3.6cm}p{5.2cm}p{5.2cm}}\toprule
& Verlet EI & WRF forward--backward\\\midrule
horizontal integrator & Störmer--Verlet, symmetric & forward--backward (conjugate to Verlet)\\
order & 2 ($\gamma=\tfrac12$) & 1 (2 in the staggered interpretation; $\beta=\tfrac12$)\\
vertical integrator & trapezoidal in $t$, $\gamma$-weighted in $W$ & $\beta$-weighted in $t$ and $W$\\
damping & none ($\gamma=\tfrac12$); $W$ off-centring & off-centring $\beta$ + divergence damping $\gamma_d$\\
unknown of the column solve & $W$ (or $\Omega$, Section~\ref{sec:omega}) & $W$ (WRF: after eliminating $\phi$; $\Omega$ diagnostic)\\
metric term & explicit (or implicit with the $\Omega$ system) & explicit\\
evaluations per small step & each part once, 1 solve (half steps combined) & each part once, 1 solve\\
stage boundaries (MIS, RK3) & no defect & first-order conjugation defect\\
\bottomrule\end{tabular}
\end{center}
WRF needs its damping: RK3 has $|g_0|=1$ for the acoustic modes (MIS study), and the forward--backward
step adds none; divergence damping and off-centring provide it (Wicker and Skamarock 2002; Baldauf
2010). With MIS methods ($|g_0|<1$) Verlet EI can be used without damping, and the off-centring $\gamma$
is available where the vertically propagating acoustic modes have to be damped.

%======================================================================
\section{Implementation in the frozen split}\label{sec:impl}
Both schemes are fast integrators of the frozen MIS split, implemented in
\code{frozen\_fast\_solve!} (\code{TimeStepping/frozen.jl}) with \code{fast = :ei} or \code{fast = :wrf}. They work for the three
frozen models \code{FVModel} (full pressure \code{:gradp} or product form), \code{ExnerModel} and
\code{MoistExnerModel}. In
the frozen fast system $Z'=F_f(C,Z)+F$ the model splits $F_f$ into the parts
\begin{center}\small
\begin{tabular}{ll}\toprule
\code{:hthermo} & $H_t$: $D_hU$ (incl.\ $-z_x\bar U$ in $\Omega$), forcing of $t$ and $W$, sponge\\
\code{:hmomentum} & $H_U$: horizontal pressure gradient $G_xp$ (incl.\ the metric term), forcing of $U$\\
\code{:vertical} & $G_v$: $D_vW$ in $t$, $G_zp-g\bar\rho$ in $W$\\\bottomrule
\end{tabular}
\end{center}
and provides the column solve $y=(I-DJ_v)^{-1}R$, $J_v$ the Jacobian of $G_v$ (exact for the linear
vertical parts of the product, Exner and moist forms, the linearly implicit approximation for
\code{:gradp}), $D$ the row scaling $c_t$ on $t$ and $c_w$ on $W$; the $W$ system depends on $c_tc_w$ only and is
factored once per stage (\code{frozen\_stage\_begin!}). Hooks (\code{FV/frozen\_ei.jl}):
\code{frozen\_vertical\_rhs!} ($DG_v$), \code{frozen\_vertical\_implicit!} ($Z\leftarrow Z+(I-DJ_v)^{-1}DG_v(Z)$) and
\code{frozen\_vertical\_solve!} (general right-hand side).

\paragraph{Options.} \code{frozen\_fast\_solve!}, \code{frozen\_mis\_step!} (MIS, Lawson, ETD) and \code{run\_frozen} take the
keywords \code{offcenter} $=\varepsilon$ and \code{smdiv} $=\gamma_d$ (default 0) for \code{:ei} and \code{:wrf} alike; \code{EI\_GAMMA[]}
(default $\tfrac12$) selects version III of EI ($\varepsilon=0$ only); \code{EI\_MERGE[]} (default true) the combined half steps.
The weights $(e_t,e_w,i_t,i_w)$ are \code{ei\_weights(h, $\varepsilon$)}.

\paragraph{Verlet EI} (\code{:ei}):
\begin{align*}
Z&\leftarrow Z+e_t\,H_t(Z)+D_{(e_t,e_w)}G_v(Z) &&E:\ t^{n+1/2},\ W\text{ explicit part}\\
Z&\leftarrow Z+h\,H_U\big(Z+\gamma_d(Z-Z_{\rm prev})\big) &&K:\ U^{n+1}\\
Z&\leftarrow Z+i_t\,H_t(Z) &&I:\ \text{horizontal part}\\
Z&\leftarrow Z+(I-D_{(i_t,i_w)}J_v)^{-1}D_{(i_t,i_w)}G_v(Z) &&I:\ t^{n+1},\ W^{n+1};
\end{align*}
$Z_{\rm prev}$ the state of the previous kick, $Z_{\rm prev}=Z$ after the first $E$. For linear $G_v$ and
$\varepsilon=0$, $\gamma=\tfrac12$ this is exactly \eqref{eq:ei1}--\eqref{eq:ei4}. Combined half steps as in Section~\ref{sec:opt}: one
horizontal momentum, one horizontal thermodynamic and one vertical evaluation and one column solve per
small step --- as WRF; only the first $E$ of a stage is extra. For \code{:gradp} (linearly implicit) and in the
sponge (evaluated at the state before $I$) combined and separate half steps differ by $O(\dtau^2)$ per step
($\le3\cdot10^{-6}$ relative in the test), and both are of order 2 for $\varepsilon=0$.

\paragraph{WRF} (\code{:wrf}):
\begin{align*}
Z&\leftarrow Z+h\,H_U\big(Z+\gamma_d(Z-Z_{\rm prev})\big) &&K:\ U^{n+1}\\
Z&\leftarrow Z+(I-\beta hJ_v)^{-1}h\big[H_t(Z)+G_v(Z)\big] &&I_{\beta h}E_{(1-\beta)h}:\ t^{n+1},\ W^{n+1},
\end{align*}
$Z_{\rm prev}=Z$ at the start of the stage. For linear $G_v$ the second line is
$Y^{n+1}=Y^n+h[H_t+(1-\beta)A_vY^n+\beta A_vY^{n+1}]$ of Section~\ref{sec:wrf}.

\paragraph{Tests} (\code{test/runtests.jl}, ``Verlet EI and WRF fast integrators''): the two-factor
column solve is exact for the linear forms; on one frozen stage of the hill problem (product,
\code{:gradp}, Exner, pressure, moist) EI is of order 2 and WRF of order 1 against RK4, both of order 1 with
$\varepsilon=0.1$; WRF$^n E=E\,$EI$^n$ for $(\varepsilon,\gamma_d)=(0,0),(0.1,0),(0.1,0.1)$; on a flat, horizontally uniform
column EI and WRF with $\varepsilon=0$ coincide with the Crank--Nicolson step of \code{:hevi} to round-off, and EI stays
bounded at the vertical Courant number 35; in MIS2 both conserve mass to round-off on terrain, EI differs
from \code{:hevi} by $O(\dtau^2)$, WRF by $O(\dtau)$; with $\gamma=0.6$ (version III) EI is of order 1. Combined and separate half steps of EI agree to round-off for the linear forms ($\gamma=0.5$
and 0.6); no small step is longer than $\Delta t/n_{\rm sub}$ (Section~\ref{sec:agnesi}). Full suite: 571 tests pass.
\begin{center}\small
\begin{tabular}{lccc}\toprule
error on one frozen stage ($T=8$\,s, max norm) & 8 steps & 16 steps & 32 steps\\\midrule
EI & 0.0242 & 0.0060 & 0.0015\\
WRF, $\beta=0.5$ & 0.116 & 0.055 & 0.027\\
WRF, $\beta=0.55$ & 0.117 & 0.057 & 0.028\\
WRF, $\beta=0.55$, $\gamma_d=0.1$ & 0.102 & 0.050 & 0.025\\\bottomrule
\end{tabular}
\end{center}
(Exner form; the other forms agree to two digits.) The linearized path (\code{fast\_solve!},
\code{:ei}, \code{:wrf} in \code{TimeStepping/fast.jl}) has the same two loops for the linear fast operator.

%======================================================================
\section{Agnesi mountain wave}\label{sec:agnesi}
Set-up as in the MIS study (\code{docs/mis\_agnesi}): hydrostatic wave over the Agnesi hill
($a=10$\,km, $h_0=400$\,m, $U=10$\,m/s, $N=0.01$\,s$^{-1}$), $200\times78$ cells, $\Delta x=2$\,km,
$\Delta\eta=200$\,m, $T=10800$\,s, $n_{\rm sub}=\lceil\Delta t/\dtau_{\rm ac}\rceil$ with $\dtau_{\rm ac}=4.3$\,s (3 small steps at
$\Delta t=10$\,s, 56 at 240\,s); reference RK4 of the same semi-discrete operator; errors are relative $L^2$
errors in $|x-x_m|\le30$\,km, $z\le9$\,km. Two forms of the extended system (\code{ExnerModel}) with an
additional pressure variable: the \emph{Exner pressure} $\pi$ ($W$ row $-c_p\Theta_f\delta\pi/\Delta z-g\bar\rho$) and the
\emph{pressure} $p$ ($W$ row $-\delta p/\Delta z-g\bar\rho$, \code{pvar = :p}); each column is in discrete hydrostatic
balance of its form (rest-state residual of the $W$ rows $4\cdot10^{-13}$). HEVI is the fast integrator of
the MIS study (horizontal Verlet half step, Crank--Nicolson vertical step, horizontal half step).
EI and WRF are run with the same options: none, $\varepsilon=0.1$, and $\varepsilon=\gamma_d=0.1$ (WRF's standard values).
\begin{center}\code{FORM=exner|pres FAST=ei|wrf|hevi [OFFCENTER=$\varepsilon$] [SMDIV=$\gamma_d$] julia examples/frozen\_agnesi.jl}\end{center}

\begin{table}[ht]\centering\footnotesize
\caption{Exner pressure: relative $L^2$ error of $\theta'$ against $\Delta t$ (s). u: unstable; bold: error $\ge0.1$.}\label{tab:theta}
\begin{tabular}{llcccccccc}\toprule
&fast integrator & 10 & 20 & 40 & 60 & 90 & 120 & 180 & 240\\\midrule
MIS2 & HEVI & 0.007 & 0.015 & 0.015 & 0.018 & 0.034 & 0.038 & \textbf{0.132} & u\\
 & EI & 0.002 & 0.002 & 0.004 & 0.005 & 0.026 & 0.034 & \textbf{0.126} & u\\
 & WRF & 0.003 & 0.002 & 0.004 & 0.005 & 0.030 & 0.044 & \textbf{0.315} & u\\
 & EI, $\varepsilon=0.1$ & 0.002 & 0.002 & 0.003 & 0.004 & 0.011 & 0.031 & \textbf{0.164} & u\\
 & WRF, $\varepsilon=0.1$ & 0.002 & 0.002 & 0.003 & 0.003 & 0.014 & 0.038 & \textbf{0.267} & u\\
 & EI, $\varepsilon=\gamma_d=0.1$ & 0.001 & 0.002 & 0.003 & 0.004 & 0.009 & 0.017 & 0.036 & u\\
 & WRF, $\varepsilon=\gamma_d=0.1$ & 0.002 & 0.002 & 0.003 & 0.003 & 0.010 & 0.020 & 0.046 & u\\
\midrule
MIS3C & HEVI & 0.005 & 0.018 & 0.035 & u & u & u & u & u\\
 & EI & 0.002 & 0.002 & 0.031 & u & u & u & u & u\\
 & WRF & 0.002 & 0.003 & 0.029 & u & u & u & u & u\\
 & EI, $\varepsilon=0.1$ & 0.002 & 0.002 & 0.002 & 0.004 & 0.006 & \textbf{0.846} & u & u\\
 & WRF, $\varepsilon=0.1$ & 0.002 & 0.002 & 0.003 & 0.005 & 0.007 & \textbf{1.340} & u & u\\
 & EI, $\varepsilon=\gamma_d=0.1$ & 0.002 & 0.002 & 0.002 & 0.004 & 0.005 & 0.008 & \textbf{0.176} & \textbf{0.261}\\
 & WRF, $\varepsilon=\gamma_d=0.1$ & 0.002 & 0.002 & 0.002 & 0.004 & 0.006 & 0.008 & \textbf{0.196} & \textbf{0.289}\\
\midrule
MIS4a & HEVI & 0.005 & 0.015 & 0.015 & 0.090 & u & \textbf{0.153} & u & u\\
 & EI & 0.001 & 0.002 & 0.004 & \textbf{0.141} & u & \textbf{1.050} & u & u\\
 & WRF & \textbf{1.186} & 0.002 & 0.004 & u & u & u & u & u\\
 & EI, $\varepsilon=0.1$ & 0.002 & 0.002 & 0.003 & 0.005 & 0.011 & 0.010 & 0.026 & \textbf{2.037}\\
 & WRF, $\varepsilon=0.1$ & \textbf{0.413} & 0.002 & 0.003 & 0.005 & 0.011 & 0.010 & 0.030 & u\\
 & EI, $\varepsilon=\gamma_d=0.1$ & 0.001 & 0.002 & 0.003 & 0.004 & 0.009 & 0.009 & 0.021 & 0.032\\
 & WRF, $\varepsilon=\gamma_d=0.1$ & 0.002 & 0.002 & 0.003 & 0.005 & 0.010 & 0.009 & 0.023 & 0.032\\
\midrule
WSRK3 & HEVI & u & u & u & u & u & u & u & u\\
 & EI & 0.013 & u & u & u & u & u & u & u\\
 & WRF & u & u & u & u & u & u & u & u\\
 & EI, $\varepsilon=0.1$ & 0.002 & 0.002 & u & u & \textbf{0.135} & u & u & \textbf{0.672}\\
 & WRF, $\varepsilon=0.1$ & 0.002 & 0.002 & u & u & \textbf{0.191} & u & u & \textbf{0.719}\\
 & EI, $\varepsilon=\gamma_d=0.1$ & 0.001 & 0.002 & u & u & 0.013 & 0.069 & \textbf{0.131} & \textbf{0.172}\\
 & WRF, $\varepsilon=\gamma_d=0.1$ & 0.002 & 0.002 & u & u & 0.013 & 0.080 & \textbf{0.145} & \textbf{0.183}\\


\bottomrule\end{tabular}
\end{table}

\begin{table}[ht]\centering\footnotesize
\caption{Pressure as additional variable: relative $L^2$ error of $\theta'$ against $\Delta t$ (s).}\label{tab:thetap}
\begin{tabular}{llcccccccc}\toprule
&fast integrator & 10 & 20 & 40 & 60 & 90 & 120 & 180 & 240\\\midrule
MIS2 & HEVI & 0.008 & 0.015 & 0.014 & 0.019 & 0.023 & 0.018 & 0.040 & u\\
 & EI & 0.003 & 0.002 & 0.006 & 0.005 & 0.019 & 0.010 & 0.033 & u\\
 & WRF & 0.003 & 0.003 & 0.006 & 0.005 & 0.019 & 0.011 & 0.032 & u\\
 & EI, $\varepsilon=0.1$ & 0.002 & 0.002 & 0.004 & 0.004 & 0.008 & 0.008 & 0.011 & u\\
 & WRF, $\varepsilon=0.1$ & 0.002 & 0.002 & 0.004 & 0.004 & 0.007 & 0.009 & 0.011 & u\\
 & EI, $\varepsilon=\gamma_d=0.1$ & 0.001 & 0.002 & 0.004 & 0.004 & 0.007 & 0.008 & 0.010 & u\\
 & WRF, $\varepsilon=\gamma_d=0.1$ & 0.002 & 0.002 & 0.004 & 0.004 & 0.007 & 0.008 & 0.010 & u\\
\midrule
MIS3C & HEVI & 0.005 & 0.018 & 0.045 & 0.027 & 0.032 & 0.061 & 0.058 & u\\
 & EI & 0.002 & 0.002 & 0.042 & 0.029 & 0.038 & 0.070 & 0.055 & 0.090\\
 & WRF & 0.002 & 0.002 & 0.039 & 0.029 & 0.037 & 0.068 & 0.054 & 0.088\\
 & EI, $\varepsilon=0.1$ & 0.002 & 0.002 & 0.002 & 0.003 & 0.003 & 0.004 & 0.006 & 0.013\\
 & WRF, $\varepsilon=0.1$ & 0.002 & 0.002 & 0.002 & 0.003 & 0.003 & 0.004 & 0.006 & 0.013\\
 & EI, $\varepsilon=\gamma_d=0.1$ & 0.002 & 0.002 & 0.002 & 0.002 & 0.003 & 0.003 & 0.005 & 0.013\\
 & WRF, $\varepsilon=\gamma_d=0.1$ & 0.002 & 0.002 & 0.002 & 0.002 & 0.003 & 0.003 & 0.004 & 0.013\\
\midrule
MIS4a & HEVI & 0.005 & 0.015 & 0.015 & 0.089 & u & 0.032 & u & u\\
 & EI & 0.001 & 0.002 & 0.003 & \textbf{0.156} & u & \textbf{0.102} & u & u\\
 & WRF & \textbf{0.500} & 0.002 & 0.004 & u & u & u & u & u\\
 & EI, $\varepsilon=0.1$ & 0.002 & 0.002 & 0.002 & 0.003 & 0.003 & 0.006 & 0.005 & 0.007\\
 & WRF, $\varepsilon=0.1$ & \textbf{0.184} & 0.002 & 0.003 & 0.003 & 0.004 & 0.007 & 0.005 & 0.007\\
 & EI, $\varepsilon=\gamma_d=0.1$ & 0.001 & 0.002 & 0.002 & 0.002 & 0.003 & 0.006 & 0.005 & 0.006\\
 & WRF, $\varepsilon=\gamma_d=0.1$ & 0.002 & 0.002 & 0.002 & 0.002 & 0.003 & 0.006 & 0.004 & 0.006\\
\midrule
WSRK3 & HEVI & u & u & 0.040 & 0.047 & 0.059 & 0.054 & 0.090 & 0.099\\
 & EI & 0.008 & 0.032 & 0.050 & 0.052 & 0.073 & 0.068 & \textbf{0.102} & \textbf{0.111}\\
 & WRF & 0.021 & 0.052 & 0.050 & 0.052 & 0.077 & 0.068 & \textbf{0.104} & \textbf{0.113}\\
 & EI, $\varepsilon=0.1$ & 0.002 & 0.002 & 0.002 & 0.004 & 0.007 & 0.005 & 0.015 & 0.016\\
 & WRF, $\varepsilon=0.1$ & 0.002 & 0.002 & 0.003 & 0.004 & 0.007 & 0.005 & 0.015 & 0.017\\
 & EI, $\varepsilon=\gamma_d=0.1$ & 0.001 & 0.002 & 0.002 & 0.003 & 0.006 & 0.004 & 0.010 & 0.016\\
 & WRF, $\varepsilon=\gamma_d=0.1$ & 0.002 & 0.002 & 0.003 & 0.004 & 0.006 & 0.004 & 0.010 & 0.016\\


\bottomrule\end{tabular}
\end{table}

\begin{table}[ht]\centering\footnotesize
\caption{Exner pressure: relative $L^2$ error of $w$. Bold: error $\ge1$.}\label{tab:w}
\begin{tabular}{llcccccccc}\toprule
&fast integrator & 10 & 20 & 40 & 60 & 90 & 120 & 180 & 240\\\midrule
MIS2 & HEVI & 0.056 & 0.104 & 0.100 & 0.107 & 0.126 & 0.205 & \textbf{1.173} & u\\
 & EI & 0.054 & 0.098 & 0.093 & 0.097 & 0.125 & 0.203 & \textbf{1.211} & u\\
 & WRF & 0.068 & 0.100 & 0.095 & 0.099 & 0.147 & 0.262 & \textbf{2.503} & u\\
 & EI, $\varepsilon=0.1$ & 0.089 & 0.095 & 0.093 & 0.096 & 0.092 & 0.191 & \textbf{1.341} & u\\
 & WRF, $\varepsilon=0.1$ & 0.089 & 0.096 & 0.094 & 0.098 & 0.106 & 0.237 & \textbf{2.088} & u\\
 & EI, $\varepsilon=\gamma_d=0.1$ & 0.091 & 0.092 & 0.092 & 0.095 & 0.087 & 0.154 & 0.350 & u\\
 & WRF, $\varepsilon=\gamma_d=0.1$ & 0.092 & 0.093 & 0.094 & 0.096 & 0.090 & 0.172 & 0.408 & u\\
\midrule
MIS3C & HEVI & 0.186 & 0.139 & 0.139 & u & u & u & u & u\\
 & EI & 0.186 & 0.136 & 0.134 & u & u & u & u & u\\
 & WRF & 0.193 & 0.150 & 0.121 & u & u & u & u & u\\
 & EI, $\varepsilon=0.1$ & 0.089 & 0.090 & 0.088 & 0.164 & 0.087 & \textbf{6.007} & u & u\\
 & WRF, $\varepsilon=0.1$ & 0.089 & 0.091 & 0.089 & 0.180 & 0.089 & \textbf{9.494} & u & u\\
 & EI, $\varepsilon=\gamma_d=0.1$ & 0.091 & 0.091 & 0.090 & 0.102 & 0.087 & 0.106 & \textbf{1.133} & \textbf{1.800}\\
 & WRF, $\varepsilon=\gamma_d=0.1$ & 0.092 & 0.092 & 0.090 & 0.104 & 0.088 & 0.107 & \textbf{1.256} & \textbf{1.972}\\
\midrule
MIS4a & HEVI & 0.154 & 0.093 & 0.097 & \textbf{7.068} & u & \textbf{1.892} & u & u\\
 & EI & 0.151 & 0.087 & 0.090 & \textbf{8.905} & u & \textbf{9.743} & u & u\\
 & WRF & \textbf{1.922} & 0.089 & 0.090 & u & u & u & u & u\\
 & EI, $\varepsilon=0.1$ & 0.090 & 0.091 & 0.091 & 0.095 & 0.150 & 0.092 & 0.173 & \textbf{14.300}\\
 & WRF, $\varepsilon=0.1$ & \textbf{8.840} & 0.093 & 0.092 & 0.097 & 0.152 & 0.092 & 0.191 & u\\
 & EI, $\varepsilon=\gamma_d=0.1$ & 0.091 & 0.091 & 0.091 & 0.096 & 0.145 & 0.086 & 0.151 & 0.085\\
 & WRF, $\varepsilon=\gamma_d=0.1$ & 0.092 & 0.092 & 0.091 & 0.096 & 0.146 & 0.086 & 0.154 & 0.090\\
\midrule
WSRK3 & HEVI & u & u & u & u & u & u & u & u\\
 & EI & \textbf{3.592} & u & u & u & u & u & u & u\\
 & WRF & u & u & u & u & u & u & u & u\\
 & EI, $\varepsilon=0.1$ & 0.093 & 0.092 & u & u & 0.934 & u & u & \textbf{2.555}\\
 & WRF, $\varepsilon=0.1$ & 0.094 & 0.093 & u & u & \textbf{1.318} & u & u & \textbf{2.780}\\
 & EI, $\varepsilon=\gamma_d=0.1$ & 0.091 & 0.091 & u & u & 0.178 & \textbf{2.278} & 0.522 & 0.589\\
 & WRF, $\varepsilon=\gamma_d=0.1$ & 0.092 & 0.092 & u & u & 0.185 & \textbf{2.526} & 0.620 & 0.631\\


\bottomrule\end{tabular}
\end{table}

\begin{table}[ht]\centering\footnotesize
\caption{Pressure as additional variable: relative $L^2$ error of $w$.}\label{tab:wp}
\begin{tabular}{llcccccccc}\toprule
&fast integrator & 10 & 20 & 40 & 60 & 90 & 120 & 180 & 240\\\midrule
MIS2 & HEVI & 0.074 & 0.099 & 0.096 & 0.104 & 0.113 & 0.111 & 0.108 & u\\
 & EI & 0.077 & 0.091 & 0.091 & 0.095 & 0.105 & 0.103 & 0.099 & u\\
 & WRF & 0.094 & 0.092 & 0.093 & 0.097 & 0.107 & 0.105 & 0.102 & u\\
 & EI, $\varepsilon=0.1$ & 0.088 & 0.095 & 0.091 & 0.095 & 0.098 & 0.102 & 0.098 & u\\
 & WRF, $\varepsilon=0.1$ & 0.088 & 0.095 & 0.092 & 0.097 & 0.100 & 0.104 & 0.101 & u\\
 & EI, $\varepsilon=\gamma_d=0.1$ & 0.091 & 0.091 & 0.091 & 0.093 & 0.097 & 0.101 & 0.096 & u\\
 & WRF, $\varepsilon=\gamma_d=0.1$ & 0.091 & 0.092 & 0.093 & 0.095 & 0.099 & 0.103 & 0.099 & u\\
\midrule
MIS3C & HEVI & 0.189 & 0.156 & 0.122 & 0.104 & 0.112 & 0.130 & 0.148 & u\\
 & EI & 0.187 & 0.149 & 0.110 & 0.096 & 0.102 & 0.125 & 0.117 & 0.141\\
 & WRF & 0.195 & 0.166 & 0.110 & 0.098 & 0.103 & 0.123 & 0.116 & 0.137\\
 & EI, $\varepsilon=0.1$ & 0.089 & 0.089 & 0.093 & 0.089 & 0.097 & 0.092 & 0.093 & 0.086\\
 & WRF, $\varepsilon=0.1$ & 0.089 & 0.089 & 0.094 & 0.090 & 0.098 & 0.093 & 0.093 & 0.087\\
 & EI, $\varepsilon=\gamma_d=0.1$ & 0.091 & 0.091 & 0.091 & 0.090 & 0.095 & 0.091 & 0.091 & 0.087\\
 & WRF, $\varepsilon=\gamma_d=0.1$ & 0.092 & 0.092 & 0.092 & 0.091 & 0.095 & 0.092 & 0.091 & 0.088\\
\midrule
MIS4a & HEVI & 0.141 & 0.106 & 0.097 & \textbf{5.841} & u & 0.883 & u & u\\
 & EI & 0.137 & 0.102 & 0.091 & \textbf{9.667} & u & \textbf{1.184} & u & u\\
 & WRF & \textbf{11.220} & 0.106 & 0.092 & u & u & u & u & u\\
 & EI, $\varepsilon=0.1$ & 0.090 & 0.092 & 0.091 & 0.089 & 0.095 & 0.098 & 0.086 & 0.094\\
 & WRF, $\varepsilon=0.1$ & \textbf{7.884} & 0.093 & 0.092 & 0.090 & 0.096 & 0.100 & 0.086 & 0.095\\
 & EI, $\varepsilon=\gamma_d=0.1$ & 0.091 & 0.091 & 0.090 & 0.090 & 0.094 & 0.097 & 0.087 & 0.093\\
 & WRF, $\varepsilon=\gamma_d=0.1$ & 0.092 & 0.092 & 0.091 & 0.090 & 0.096 & 0.099 & 0.087 & 0.094\\
\midrule
WSRK3 & HEVI & u & u & 0.143 & 0.159 & 0.204 & 0.128 & 0.133 & 0.162\\
 & EI & 0.174 & 0.149 & 0.146 & 0.139 & 0.223 & 0.098 & 0.154 & 0.140\\
 & WRF & \textbf{7.230} & \textbf{4.551} & 0.144 & 0.145 & 0.234 & 0.114 & 0.159 & 0.141\\
 & EI, $\varepsilon=0.1$ & 0.093 & 0.094 & 0.089 & 0.087 & 0.087 & 0.080 & 0.092 & 0.080\\
 & WRF, $\varepsilon=0.1$ & 0.094 & 0.095 & 0.090 & 0.088 & 0.087 & 0.080 & 0.092 & 0.077\\
 & EI, $\varepsilon=\gamma_d=0.1$ & 0.091 & 0.090 & 0.089 & 0.087 & 0.082 & 0.082 & 0.085 & 0.083\\
 & WRF, $\varepsilon=\gamma_d=0.1$ & 0.092 & 0.091 & 0.090 & 0.089 & 0.082 & 0.081 & 0.085 & 0.082\\


\bottomrule\end{tabular}
\end{table}

\paragraph{Number of small steps per stage (corrected).} Stage $i$ of a MIS step integrates over $d_i\Delta t$.
The number of small steps was $\max(1,\mathrm{round}(d_in_{\rm sub}))$, which allows small steps up to
$1.5\,\Delta t/n_{\rm sub}$, i.e.\ above the horizontal acoustic limit for which $n_{\rm sub}$ was chosen. MIS2
($d=(0.127,0.590,0.470)$) with $n_{\rm sub}=3$ at $\Delta t=10$\,s ran its third stage in one small step of
$4.7\,$s$\,>\dtau_{\rm ac}=4.3$\,s, and with $n_{\rm sub}=5$ at 20\,s in two steps of 4.7\,s: exactly the two cases in
which MIS2 with EI and WRF was unstable in the first version of this note. HEVI survived them, because its
horizontal Verlet steps are half steps. The count is now
\[
n_i=\max\big(1,\lceil d_in_{\rm sub}\rceil\big)\qquad(\code{fast\_substeps}),
\]
so that no small step is longer than $\Delta t/n_{\rm sub}$, for MIS, ETD, Lawson and EPIRK and in the frozen and
the linearized split. The MIS study of \code{docs/mis\_agnesi} used the rounded count; its HEVI
numbers change slightly (e.g.\ MIS2 at 10\,s: $\theta'$ error 0.014 $\to$ 0.007). All results below use
the corrected count.

\paragraph{EI and WRF with the same options behave the same.} In both forms and for all three option
sets the stable and unstable cases of EI and WRF coincide, and the errors agree to one or two digits,
EI's being slightly smaller (no change of variables missing at the stage boundaries). The few
differences are at the edges, where the stage-boundary defect of WRF matters: MIS4a at 10\,s
($n_{\rm sub}=3$, first stage one step) is weakly unstable with WRF for $\gamma_d=0$ (Exner: $\theta'$ error 1.19 and 0.41;
pressure: 0.50 and 0.18) and fine with EI; WSRK3 at 10 and 20\,s without options ($w$ error 7.2 and 4.6 with
WRF against 0.17 and 0.15 with EI, pressure form).

\paragraph{Accuracy.} With $\varepsilon=0.1$, EI and WRF reduce the $\theta'$ error at $\Delta t\le60$\,s from
$0.5$--$4.5\cdot10^{-2}$ (HEVI) to $1.4$--$5\cdot10^{-3}$, for all four methods and both forms (without options at
$\Delta t\le20$\,s): with HEVI the $\theta'$
error is dominated by the horizontal--vertical splitting of the fast step (the Strang splitting of the
stiff vertical part), not by the MIS method. EI and WRF use $D_hU^{n+1}$ and $D_vW^{n+1}$ in the same
thermodynamic update and avoid it. The $w$ error ($\approx0.09$) is the same for all stable runs: the undamped
vertical acoustic modes of the MIS study.

\paragraph{Off-centring and extrapolation.}
\begin{itemize}
\item Without options EI, WRF and HEVI have irregular stability gaps in the Exner form (MIS3C from 60\,s,
MIS4a at 60--120\,s, WSRK3 almost everywhere); in the pressure form MIS2 (up to 180\,s), MIS3C and
WSRK3 are stable (HEVI: WSRK3 unstable at 10--20\,s), MIS4a still has the gap at 60--90\,s and from 180\,s
(WRF from 60\,s).
\item $\varepsilon=0.1$ removes the gaps at small and medium $\Delta t$ and reduces the errors at $\Delta t\ge40$\,s by a factor 2--20. In
the pressure form all four methods are then stable up to 240\,s (MIS2 up to 180\,s), WSRK3 included,
with $\theta'$ errors $\le0.017$.
\item $\gamma_d=0.1$ matters only in the Exner form at $\Delta t\ge120$\,s (MIS3C, MIS4a and WSRK3 stable up to 240\,s
except WSRK3 at 40--60\,s); in the pressure form it changes the errors by a few per cent.
\end{itemize}

\paragraph{Pressure against Exner pressure.} With the pressure as additional variable all methods are
markedly more robust: with $\varepsilon=0.1$, MIS3C at 240\,s has $\theta'$ error 0.013 (Exner: unstable; 0.26 with
$\gamma_d=0.1$), MIS4a 0.0065 (Exner: 2.0 or unstable; 0.032 with $\gamma_d$), WSRK3 0.016 (Exner with $\gamma_d$: unstable at
40--60\,s, 0.17 at 240\,s), and HEVI makes
WSRK3 stable from 40\,s. A plausible cause, not analysed here: in the pressure form the $W$ row $-\delta p$ has
no frozen coefficient, while the Exner form carries the frozen $c_p\Theta_f$ in front of $\delta\pi$ and in the
$\pi$ equation, so that the frozen fast operator of a stage and the slow remainder differ more (the
constraint-defect modes of the MIS study).

\paragraph{Cost.} Serial runs, MIS4a at $\Delta t=60$\,s ($n_{\rm sub}=14$, $T=10800$\,s, the first $\Delta t$ of each
run excluded because of compilation; Exner form): HEVI 89\,s, EI (combined half steps, version III $\gamma=0.55$)
49\,s, WRF ($\varepsilon=\gamma_d=0.1$) 44\,s. The combined EI costs 10\,\% more than WRF (the first $E$ of each stage and, for
version III, the column-local correction) and 55\,\% of HEVI.

%======================================================================
\section{Open decisions}
\begin{enumerate}
\item Default fast integrator: EI or WRF, both with $\varepsilon=0.1$ (and $\gamma_d=0.1$ for the Exner form). They are
conjugate; EI avoids the stage-boundary defect at the same cost.
\item Default pressure variable of the extended system: $p$ is markedly more robust than $\pi$ on
Agnesi; to be checked on the moist cases (\code{MoistExnerModel}, \code{pvar = :p}).
\item The remaining gaps (Exner form; MIS4a without options): linear stability analysis of the coupled
MIS--EI/WRF step.
\item Rerun of the MIS study (\code{docs/mis\_agnesi}) with the corrected number of small steps.
\item $W$ system with explicit metric term (sufficient for $c_s\dtau|z_x|/(J\Delta\eta)\lesssim1$) or the $\Omega$
system with implicit metric term; for the $\Omega$ system the column-local approximations (a), (b) or the exact
two-dimensional system.
\end{enumerate}

%======================================================================
\begin{thebibliography}{9}
\bibitem{bal} M.~Baldauf: Linear stability analysis of Runge--Kutta-based partial time-splitting
schemes for the Euler equations. Mon.\ Wea.\ Rev.\ 138 (2010) 4475--4496.
\bibitem{gal} T.~Gal-Chen, R.~C.~J.~Somerville: On the use of a coordinate transformation for the
solution of the Navier--Stokes equations. J.\ Comput.\ Phys.\ 17 (1975) 209--228.
\bibitem{hlw} E.~Hairer, C.~Lubich, G.~Wanner: Geometric Numerical Integration, 2nd ed. Springer
2006.
\bibitem{imex} Split-explicit with implicit treatment of part of the fast waves (note, 4 June 2024;
Imex.pdf).
\bibitem{ksd} J.~B.~Klemp, W.~C.~Skamarock, J.~Dudhia: Conservative split-explicit time integration
methods for the compressible nonhydrostatic equations. Mon.\ Wea.\ Rev.\ 135 (2007) 2897--2913.
\bibitem{wrf} W.~C.~Skamarock et al.: A description of the Advanced Research WRF Model Version 4.
NCAR Technical Note NCAR/TN-556+STR (2019).
\bibitem{ws} L.~J.~Wicker, W.~C.~Skamarock: Time-splitting methods for elastic models using forward
time schemes. Mon.\ Wea.\ Rev.\ 130 (2002) 2088--2097.
\end{thebibliography}
\end{document}
